% Generated draft: requires mathematical and source review.
\documentclass[11pt]{article}
\usepackage[margin=25mm]{geometry}
\usepackage{fontspec}
\setmainfont{Latin Modern Roman}
\usepackage{amsmath,amssymb,mathtools,unicode-math}
\setmathfont{Latin Modern Math}
\usepackage{enumitem,needspace,xurl,hyperref}
\setlength{\parindent}{0pt}
\setlength{\parskip}{5pt}
\setlength{\emergencystretch}{3em}
\newcommand{\N}{\mathbb{N}}
\newcommand{\Z}{\mathbb{Z}}
\newcommand{\Q}{\mathbb{Q}}
\newcommand{\R}{\mathbb{R}}
\newcommand{\C}{\mathbb{C}}
\newcommand{\F}{\mathbb{F}}
\newcommand{\A}{\mathbb{A}}
\newcommand{\PP}{\mathbb{P}}
\newcommand{\E}{\mathbb{E}}
\newcommand{\eps}{\varepsilon}
\newcommand{\dd}{\,\mathrm{d}}
\begin{document}
\section*{UnsolvedMath ID: 30004755}
\textbf{Source problem code:} catalogue entry
\subsection{Definitions and mathematical statement}
Let \(X\) be a closed, oriented, smooth four-manifold. Write
\[
b_1(X)=\dim_{\R}H^1(X;\R),
\]
and let \(b^+(X)\) be the positive index of the intersection form on \(H^2(X;\R)\). Let \(\chi(X)\) denote the Euler characteristic and \(\sigma(X)\) the signature of this form.

For the proposed inequality, interpret the Chern-number notation by the standard topological convention
\[
c_2(X):=\chi(X),\qquad c_1(X)^2:=2\chi(X)+3\sigma(X).
\]
The supplied record does not explicitly define these symbols. This convention specifies numerical invariants; it does not assume that \(X\) has a chosen complex structure or a distinguished first Chern class.

For a \(\operatorname{Spin}^c\) structure \(\mathfrak{s}\), write
\[
c_1(\mathfrak{s})^2
:=\left\langle c_1(\mathfrak{s})\smile c_1(\mathfrak{s}),[X]\right\rangle,
\]
and write \(\operatorname{SW}_X(\mathfrak{s})\) for its numerical Seiberg--Witten invariant. In the usual connected-manifold convention, Seiberg--Witten simple type means that, for every \(\mathfrak{s}\),
\[
\operatorname{SW}_X(\mathfrak{s})\ne0
\quad\Longrightarrow\quad
c_1(\mathfrak{s})^2=2\chi(X)+3\sigma(X).
\]
The nonvanishing hypothesis means that there exists at least one \(\mathfrak{s}\) with \(\operatorname{SW}_X(\mathfrak{s})\ne0\). A homology orientation is understood when defining the signed invariants; their nonvanishing does not depend on that sign choice. The record does not explicitly specify connectedness or a convention for Seiberg--Witten invariants of disconnected manifolds, so that scope qualification remains unresolved.

\textbf{Supplied conjectural statement.} For every closed, oriented, smooth four-manifold \(X\) satisfying
\[
b_1(X)=0,\qquad b^+(X)\ge3,\qquad b^+(X)\equiv1\pmod2,
\]
which has Seiberg--Witten simple type and a nonzero Seiberg--Witten invariant, must
\[
c_1(X)^2\le3c_2(X)
\]
hold? With the numerical convention above, the conclusion is equivalently
\[
2\chi(X)+3\sigma(X)\le3\chi(X),
\qquad\text{or}\qquad
3\sigma(X)\le\chi(X).
\]
The interpretation of the Chern-number notation and the connectedness scope require confirmation against the source. The supplied registry labels the problem open; this draft does not independently establish its literature status.
\subsection{Short English statement}
The supplied problem asks whether every closed, oriented, smooth four-manifold with vanishing first Betti number, odd positive intersection-form index at least three, Seiberg--Witten simple type, and a nonzero Seiberg--Witten invariant satisfies a Bogomolov--Miyaoka--Yau type inequality. Under the stated topological Chern-number convention, this asks whether \(3\sigma(X)\le\chi(X)\). The notation and connectedness convention require source confirmation.
\subsection{Sources}
\begin{itemize}
\item Supplied local UnsolvedMath registry record, permanent ID \(30004755\), source problem code \texttt{catalogue entry}, titled ``Bogomolov--Miyaoka--Yau Inequality for Smooth Four-Manifolds,'' in the Oberwolfach Reports Open Problems collection. No individual proposer is supplied; the recorded proposed year is 2021.
\item Supplied source citation: \emph{Partial Differential Equations (2022)}, Oberwolfach Reports, pp.\ 1859--1914, DOI 10.4171/owr/2021/35. \url{https://doi.org/10.4171/owr/2021/35}. The supplied background describes Oberwolfach Reports as documenting mathematical workshops and their open-problem discussions.
\item Supplied literature source, labeled ``Feehan--Leness'': \url{https://arxiv.org/abs/2010.15789}.
\item Supplied literature source, labeled ``Feehan--Leness sequel'': \url{https://arxiv.org/abs/2410.13809}.
\item Supplied literature source, labeled ``Oberwolfach report'': \url{https://ems.press/content/serial-article-files/46911}.
\item Supplied literature assessment, checked 2026-08-22: the specified inequality is labeled open and conjectural, with later work described as developing the program without proving it in full. This is supplied dated curation, not an independently verified assessment; the record explicitly requests independent verification of time-sensitive status and cited claims.
\item The registry marks the statement \texttt{corrected\_verified} and reports that the official report was checked directly, with definitions restored and extraction artifacts or source-dependent wording removed. This is a supplied verification claim; no source visit or independent verification is asserted here.
\end{itemize}
\end{document}
